\section*{अध्याय \ref{ch:b1:rock-forming-minerals} --- मुख्य शैल-निर्माता खनिज}

\begin{solution}{exo:b1:rock-forming-minerals:1}
\ce{SiO4^{4-}}: \omterm{def:b1:rock-forming-minerals:silicate-classes}{नीसोसिलिकेट}, $s = 0$। \ce{Si2O7^{6-}}: \omterm{def:b1:rock-forming-minerals:silicate-classes}{सोरोसिलिकेट},
$s = 1$। \ce{Si6O18^{12-}}: \omterm{def:b1:rock-forming-minerals:silicate-classes}{साइक्लोसिलिकेट}, $s = 2$। \ce{Si2O5^{2-}}:
\omterm{def:b1:rock-forming-minerals:silicate-classes}{फ़िलोसिलिकेट}, $s = 3$। \ce{SiO2}: \omterm{def:b1:rock-forming-minerals:silicate-classes}{टेक्टोसिलिकेट}, $s = 4$। \ce{SiO3^{2-}}
के लिए $4 - s/2 = 3$ है, अतः $s = 2$: एक वलय (\omterm{def:b1:rock-forming-minerals:silicate-classes}{साइक्लोसिलिकेट}) या एक एकल शृंखला
(\omterm{def:b1:rock-forming-minerals:silicate-classes}{इनोसिलिकेट}); अनुपात अकेले उन्हें अलग नहीं कर सकता।
\end{solution}

\begin{solution}{exo:b1:rock-forming-minerals:2}
फ़ॉर्स्टेराइट: $2\times24.305 + 28.085 + 4\times15.999 =
\qty{140.691}{g.mol^{-1}}$; फ़ायलाइट: $2\times55.845 + 28.085 +
4\times15.999 = \qty{203.771}{g.mol^{-1}}$। मोलर आयतन:
$140.691/3.275 = \qty{42.96}{cm^3.mol^{-1}}$ और $203.771/4.392 =
\qty{46.40}{cm^3.mol^{-1}}$।
\end{solution}

\begin{solution}{exo:b1:rock-forming-minerals:3}
$(0.78 - 0.72)/0.72 = +\qty{8.3}{\percent}$ \ce{Fe^{2+}} के लिए;
$-\qty{4.2}{\percent}$ \ce{Ni^{2+}} के लिए; $+\qty{15.3}{\percent}$ \ce{Mn^{2+}} के लिए; $+\qty{38.9}{\percent}$ \ce{Ca^{2+}}
के लिए। लोहा और निकेल, \qty{15}{\percent} से कम, मैग्नीशियम को मुक्त रूप से बदलते हैं; मैंगनीज़, ठीक ऊपर, सीमित है;
कैल्शियम, \qty{30}{\percent} से आगे, मुश्किल से प्रवेश करता है। चारों का आवेश मैग्नीशियम के समान है।
\end{solution}

\begin{solution}{exo:b1:rock-forming-minerals:4}
$(3.51 - 2.16)/2.16 = \qty{62.5}{\percent}$ ग्रेफ़ाइट से हीरे तक;
$(2.93 - 2.71)/2.71 = \qty{8.1}{\percent}$ कैल्साइट से ऐरागोनाइट तक। घने \omterm{def:b1:rock-forming-minerals:polymorph}{बहुरूप}, हीरा और ऐरागोनाइट, उच्च-दाब रूप हैं।
\end{solution}

\begin{solution}{exo:b1:rock-forming-minerals:5}
$x = 0.30$ के साथ $\mathrm{Na}_{1-x}\mathrm{Ca}_x\mathrm{Al}_{1+x}\mathrm{Si}_{3-x}\mathrm{O}_8$ में:
$\mathrm{(Na_{0.7}Ca_{0.3})(Al_{1.3}Si_{2.7})O_8}$। आवेश:
$0.7 + 2\times0.3 + 3\times1.3 + 4\times2.7 = 0.7 + 0.6 + 3.9 + 10.8 = +16$,
और $8\times(-2) = -16$।
\end{solution}

\begin{solution}{exo:b1:rock-forming-minerals:6}
एक दोहराते समूह में दो चतुष्फलक तीन कोने साझा करते हैं (वे जिनका शीर्ष दोनों शृंखलाओं को जोड़ता है) और दो दो। चार
चतुष्फलकों के $4\times4 = 16$ कोनों में से $3 + 3 + 2 + 2 = 10$ साझा हैं, अर्थात्
$10/2 = 5$ साझा ऑक्सीजन; शेष $16 - 10 = 6$ अ-साझा हैं। चार सिलिकनों के लिए कुल
$5 + 6 = 11$ ऑक्सीजन: \ce{Si4O11}, आवेश
$4\times4 - 11\times2 = -6$ के साथ। प्रस्ताव $s = 5/2$ के साथ वही देता है: प्रति सिलिकन
$4 - 5/4 = 11/4$ ऑक्सीजन, प्रति सिलिकन आवेश $5/2 - 4 = -3/2$।
\end{solution}

\begin{solution}{exo:b1:rock-forming-minerals:7}
\qty{100}{g} में मोल: \ce{SiO2} $40.20/60.083 = 0.66907$; \ce{MgO}
$47.10/40.304 = 1.16862$; \ce{FeO} $12.20/71.844 = 0.16981$; \ce{NiO}
$0.40/74.692 = 0.00536$। ऑक्सीजन: $2\times0.66907 + 1.16862 + 0.16981 +
0.00536 = 2.68194$; $f = 4/2.68194 = 1.49146$। प्रति चार ऑक्सीजन: Si
\num{0.998}, Mg \num{1.743}, Fe \num{0.253}, Ni \num{0.008}। अष्टफलकीय स्थल
$1.743 + 0.253 + 0.008 =$ $2.004$ रखते हैं, और चतुष्फलकीय स्थल
\num{0.998}: एक अच्छा विश्लेषण, \ce{(Mg_{1.743}Fe_{0.253}Ni_{0.008})SiO4}।
$\mathrm{Fo} = 100\times1.743/(1.743 + 0.253) = 87.3$। गुणक $f$ Mg और Fe को एक-सा गुणा करता है, अतः अनुपात में
वह निरस्त हो जाता है, अर्थात्:
$\mathrm{Fo} =$ $100\,(47.10/40.304)/(47.10/40.304 + 12.20/71.844) =$
$100\times1.16862/1.33843 = 87.3$।
\end{solution}

\begin{solution}{exo:b1:rock-forming-minerals:8}
धनात्मक: $1 + 2\times3 + 3\times4 + 3 = +22$; ऋणात्मक:
$10\times(-2) + 2\times(-1) = -22$। \ce{Si4O10} संरचना की एक परत पर $-4$ आवेश होता
(\cref{prop:b1:rock-forming-minerals:si-o}, $s = 3$); चार में से एक सिलिकन को एल्युमिनियम से बदलने पर एक ऋणात्मक आवेश
जुड़ता है, अतः \ce{(Si3Al)O10} पर $-5$ होता है। दो अष्टफलकीय एल्युमिनियम दो हाइड्रॉक्सिलों के साथ $6 - 2 = +4$ देते
हैं, जो पाँच में से चार को संतुलित करते हैं; परतों के बीच का पोटैशियम आयन पाँचवें को संतुलित करता है।
\end{solution}

\begin{solution}{exo:b1:rock-forming-minerals:9}
(क) $M = (140.691 + 203.771)/2 = \qty{172.231}{g.mol^{-1}}$, $V = (42.96 +
46.40)/2 = \qty{44.68}{cm^3.mol^{-1}}$, $\rho = 172.231/44.68 =
\qty{3.855}{g.cm^{-3}}$; घनत्वों का माध्य $(3.275 + 4.392)/2 =
\qty{3.834}{g.cm^{-3}}$ है, \qty{0.021}{g.cm^{-3}} कम। (ख) हल कीजिए
$203.771 + x(140.691 - 203.771) = 3.600\,\bigl(46.40 + x(42.96 - 46.40)\bigr)$:
$203.771 - 63.080x = 167.040 - 12.384x$, अतः $x = 36.731/50.696 = 0.7245$: \omterm{def:b1:rock-forming-minerals:groups}{ओलिवीन}
$\text{Fo}_{72}$ है (निकटतम प्रतिशत तक)।
\end{solution}

\begin{solution}{exo:b1:rock-forming-minerals:10}
$M_i = \rho_i V_i$ के साथ,
\[
  \rho(x) - \bigl(x\rho_1 + (1 - x)\rho_2\bigr)
  = \frac{x\rho_1V_1 + (1 - x)\rho_2V_2 - \bigl(x\rho_1 + (1 - x)\rho_2\bigr)\bigl(xV_1 + (1 - x)V_2\bigr)}{V(x)} .
\]
गुणनफल का विस्तार करने पर अंश है $x\rho_1V_1 - x^2\rho_1V_1 +
(1 - x)\rho_2V_2 - (1 - x)^2\rho_2V_2 - x(1 - x)(\rho_1V_2 + \rho_2V_1)$, और
$x - x^2 = x(1 - x)$, $(1 - x) - (1 - x)^2 = x(1 - x)$ है, अतः वह
$x(1 - x)(\rho_1V_1 + \rho_2V_2 - \rho_1V_2 - \rho_2V_1) =
x(1 - x)(\rho_1 - \rho_2)(V_1 - V_2)$ है। \omterm{def:b1:rock-forming-minerals:groups}{ओलिवीन} के लिए (1 = फ़ॉर्स्टेराइट) दोनों अंतर ऋणात्मक हैं,
$3.275 - 4.392 = -1.117$ और $42.96 - 46.40 =
-3.44$, अतः उनका गुणनफल धनात्मक है और घनत्व सीधी रेखा से ऊपर है। $\text{Fo}_{50}$ पर:
$0.25\times1.117\times3.44/44.68 =
\qty{0.0215}{g.cm^{-3}}$, जो \cref{exo:b1:rock-forming-minerals:9} से मोलर आयतनों के पूर्णांकन तक मेल खाता है।
\end{solution}

\begin{solution}{exo:b1:rock-forming-minerals:11}
कुल \qty{99.60}{\percent}। ऑक्साइडों के मोल: \ce{SiO2} $0.90042$,
\ce{Al2O3} $0.28344$, \ce{CaO} $0.19973$, \ce{Na2O} $0.08229$,
\ce{K2O} $0.00318$। ऑक्सीजन: $1.80084 + 0.85033 + 0.19973 + 0.08229 +
0.00318 = 2.93636$; $f = 8/2.93636 = 2.72446$। धनायन (\ce{Al2O3}, \ce{Na2O}, \ce{K2O} के लिए मोलों का दुगुना): Si \num{2.453}, Al \num{1.544}, Ca
\num{0.544}, Na \num{0.448}, K \num{0.017}। चतुष्फलकीय स्थल:
$2.453 + 1.544 = 3.997$; बड़ा स्थल: $0.544 + 0.448 + 0.017 = 1.009$। तब
$\mathrm{Al} - 1 = 0.544$, Ca के बराबर पूर्णांकन की सीमा तक: हर कैल्शियम एक एल्युमिनियम के साथ आया था। $\mathrm{An} = 100\times0.544/1.009 = 53.9$: एक
लैब्रेडोराइट, $\text{An}_{54}$।
\end{solution}

\begin{solution}{exo:b1:rock-forming-minerals:12}
(क) $M(\ce{Fe3O4}) = 3\times55.845 + 4\times15.999 = \qty{231.531}{g.mol^{-1}}$।
उसके तीन लोहे, \ce{FeO} के रूप में दिए गए, $3\times71.844 = \qty{215.532}{g}$ भारी हैं,
यानी $215.532/231.531 = \qty{93.09}{\percent}$; \ce{Fe^{3+}} से बँधी ऑक्सीजन, जो \ce{FeO} के हिसाब से परे है, योग में
नहीं है। (ख) $93.09/71.844 = 1.29572$ मोल \ce{FeO} और ऑक्सीजन के; $f =
4/1.29572 = 3.08708$; चार ऑक्सीजन पर चार लोहे, \ce{Fe4O4}। मैग्नेटाइट में चार ऑक्सीजन पर तीन धनायन हैं। ऑक्सीजन कम
गिनी गईं: हर \ce{Fe^{3+}} डेढ़ ऑक्सीजन रखता है, एक नहीं जैसा \ce{FeO} मानता है। (ग) तीन धनायनों पर, Fe $= 3$
के साथ $a$ \ce{Fe^{2+}} और $b$ \ce{Fe^{3+}}: $a + b = 3$ और $2a + 3b = 8$,
अतः $b = 2$, $a = 1$: $\mathrm{Fe^{2+}Fe^{3+}_2O_4}$। (घ) प्रति सूत्र-एकक एक \ce{FeO} और
एक \ce{Fe2O3}: $71.844/231.531 = \qty{31.03}{\percent}$
\ce{FeO} और $159.687/231.531 = \qty{68.97}{\percent}$ \ce{Fe2O3}, कुल
\qty{100.00}{\percent}।
\end{solution}

\begin{solution}{pb:b1:rock-forming-minerals:1}
\textbf{1.} \ce{SiO2} $28.085 + 2\times15.999 = 60.083$; \ce{FeO} 71.844;
\ce{MnO} 70.937; \ce{MgO} 40.304; \ce{NiO} 74.692; \ce{CaO} 56.077
(\si{g.mol^{-1}})।

\textbf{2.} $39.50 + 16.53 + 0.23 + 43.32 + 0.20 + 0.22 =
\qty{100.00}{\percent}$। \qty{97}{\percent} का योग एक ऐसे घटक का संकेत देता जो विश्लेषित नहीं हुआ (पानी, या
\ce{Fe^{3+}} से बँधी ऑक्सीजन जो \ce{FeO} के रूप में दी गई), या एक खराब विश्लेषण।

\textbf{3.} \ce{SiO2} $39.50/60.083 = 0.65742$; \ce{FeO} 0.23008; \ce{MnO}
0.00324; \ce{MgO} 1.07483; \ce{NiO} 0.00268; \ce{CaO} 0.00392।

\textbf{4.} \ce{SiO2} के लिए $2\times0.65742 = 1.31485$, बाकी के लिए प्रश्न 3 के समान संख्याएँ; योग $1.31485 + 0.23008 + 0.00324 + 1.07483 +
0.00268 + 0.00392 = 2.62960$।

\textbf{5.} हर ऑक्साइड में एक धनायन है: प्रश्न 3 की संख्याएँ।

\textbf{6.} $f = 4/2.62960 = 1.52114$।

\textbf{7.} Si $0.65742\times1.52114 = 1.000$; Fe 0.350; Mn 0.005; Mg 1.635;
Ni 0.004; Ca 0.006।

\textbf{8.} अष्टफलकीय स्थलों के लिए $1.635 + 0.350 + 0.005 + 0.004 + 0.006 = 2.000$ और चतुष्फलकीय स्थल के लिए 1.000: ठीक
आदर्श 2 और 1।

\textbf{9.} \ce{(Mg_{1.635}Fe_{0.350}Mn_{0.005}Ni_{0.004}Ca_{0.006})SiO4}।

\textbf{10.} $2\times2.000 + 4\times1.000 = +8.000$, $4\times(-2) =
-8$ के सामने। हर धनायन को उसके ऑक्साइड के आवेश से गिना गया था, इसलिए धनात्मक आवेश संरचना से ही ऑक्सीजन के दुगुने
हैं: जाँच विफल नहीं हो सकती थी। प्रश्न 8 के स्थल-योग विफल हो सकते थे।

\textbf{11.} निकेल मैग्नीशियम से
$(0.69 - 0.72)/0.72 = -\qty{4}{\percent}$ से भिन्न है, \qty{15}{\percent} के भीतर ही: वह मुक्त रूप से
प्रतिस्थापित होता है। कैल्शियम $(1.00 - 0.72)/0.72 = \qty{39}{\percent}$ बड़ा है, \qty{30}{\percent} से आगे: वह मुश्किल से प्रवेश करता है।

\textbf{12.} लोहा जब \ce{FeO} के रूप में गिना जाए और उसका एक भाग \ce{Fe^{3+}} हो, तो प्रति धनायन कम ऑक्सीजन
मिलती है, अतः चार ऑक्सीजन पर तीन से अधिक धनायन
(\cref{exo:b1:rock-forming-minerals:12})। योग $1.000 + 2.000 = 3.000$ है: कोई महत्वपूर्ण \ce{Fe^{3+}} नहीं।

\textbf{13.} $\mathrm{Fo} = 100\times1.635/(1.635 + 0.350) = 163.5/1.985 =
82.4$।

\textbf{14.} $\mathrm{Mg\#} = 82.4$: सारा लोहा \ce{Fe^{2+}} है (प्रश्न 12), अतः $\mathrm{Fe^{2+}} = \mathrm{Fe}$ और दोनों
परिभाषाएँ एक हो जाती हैं।

\textbf{15.} $\text{Fo}_{82.4}\text{Fa}_{17.6}$, $x = 0.824$; द्विआधारी प्रतिरूप में लोहा अष्टफलकीय स्थलों के
\qty{17.6}{\percent} भरता है ($0.350/2.000 = \qty{17.5}{\percent}$ लघु तत्वों के साथ)।

\textbf{16.} गुणक Mg और Fe को एक-सा गुणा करता है और निरस्त हो जाता है:
$\mathrm{Fo} = 100\times1.07483/(1.07483 + 0.23008) = 82.4$।

\textbf{17.} $40.5/40.304 = 1.00486$ और $20.0/71.844 = 0.27838$;
$\mathrm{Fo} = 100\times1.00486/1.28324 = 78.3$। बाद का \omterm{def:b1:rock-forming-minerals:groups}{ओलिवीन} लोहे में अधिक समृद्ध है।

\textbf{18.} $M = 140.691$ और \qty{203.771}{g.mol^{-1}};
$V = 140.691/3.275 = 42.96$ और $203.771/4.392 =
\qty{46.40}{cm^3.mol^{-1}}$।

\textbf{19.} $M = 0.824\times140.691 + 0.176\times203.771 = 115.929 +
35.864 = \qty{151.793}{g.mol^{-1}}$; $V = 0.824\times42.96 +
0.176\times46.40 = 35.399 + 8.166 = \qty{43.565}{cm^3.mol^{-1}}$।

\textbf{20.} $\rho = 151.793/43.565 = \qty{3.484}{g.cm^{-3}}$।

\textbf{21.} $0.824\times3.275 + 0.176\times4.392 = 2.699 + 0.773 =
\qty{3.472}{g.cm^{-3}}$, आदर्श-मिश्रण मान से \qty{0.012}{g.cm^{-3}} (\qty{0.3}{\percent}) नीचे: निकट, जैसा
\cref{exo:b1:rock-forming-minerals:10} पूर्वानुमान करता है।

\textbf{22.} $0.12\times\qty{3484}{kg.m^{-3}}\times\qty{1}{m^3} =
\qty{418}{kg}$।

\textbf{23.} $418/2950 = \qty{14.2}{\percent}$: सघन \omterm{def:b1:rock-forming-minerals:groups}{ओलिवीन} आयतन की तुलना में द्रव्यमान का एक बड़ा भाग बनाता है।

\textbf{24.} \omterm{def:b1:rock-forming-minerals:groups}{ओलिवीन} $\text{Fo}_{82}$ है: उसकी फ़ॉर्स्टेराइट मात्रा \textbf{82 प्रतिशत} है।
\end{solution}
